%Chargement des paquets
\usepackage{amsmath}
\usepackage{amsfonts}
\usepackage{amssymb}
\usepackage{enumerate}
\usepackage{mathtools}
\usepackage{bbm}
\usepackage{xparse, etoolbox}
\usepackage{enumerate}
\usepackage{enumitem}
\usepackage{mathabx}
\usepackage{babel}[french]

%environment
\usepackage{amsthm}
\usepackage{keytheorems}
\usepackage{hyperref} 

%Interface théorème
\newkeytheorem{lem}[
	name = Lemme
]

\newkeytheorem{prop}[
	name = Proposition
]

\newkeytheorem{thm}[
	name = Théorème
]

\newkeytheorem{coro}[
	name = Corollaire
]

\renewcommand*{\proofname}{Démonstration}

\theoremstyle{definition}
\newtheorem{defn}{Définition}

\theoremstyle{definition}
\newtheorem*{nota}{Notation}

\theoremstyle{definition}
\newtheorem*{limi}{Liminaire}

\theoremstyle{remark}
\newtheorem{rmq}{Remarque}

\theoremstyle{plain}
\newtheorem*{fait}{Fait}


%Ensembles de nombres
\newcommand{\N} {\mathbb{N}}
\newcommand{\Ne}{\N^\ast}
\newcommand{\Z} {\mathbb{Z}}
\newcommand{\Q} {\mathbb{Q}}
\newcommand{\R} {\mathbb{R}}
\newcommand{\Rb}{\overline{\mathbb{R}}}
\newcommand{\Rp}{\R_+}
\newcommand{\Rm}{\R_-}
\newcommand{\K} {\mathbb{K}}
\newcommand{\Ket} {\mathbb{K^{\ast}}}
\newcommand{\Cx}{\mathbb{C}}

%Ensembles, tribus
\newcommand{\A}{\mathbb{A}}
\renewcommand{\P}{\mathcal{P}}
\newcommand{\T}{\mathcal{T}}
\newcommand{\F}{\mathcal{F}}
\newcommand{\C} {\mathcal{C}}
\newcommand{\ouv}{\mathcal{O}}
\newcommand{\B}{\mathcal{B}}
\newcommand{\M}{\mathcal{M}}
\newcommand{\etag}{\mathcal{E}}
\newcommand{\etagOp}{\etag^{+}(\Omega)}
\newcommand{\etagO}{\etag(\Omega)}
%%Lp avant quotientage
\newcommand{\Lic}{\mathcal{L}^1}
\newcommand{\Lpc}{\mathcal{L}^p}
\newcommand{\Linfc}{\mathcal{L}^{\infty}}
\newcommand{\Lifull}{\Lic(\Omega, \B, m)}
\newcommand{\Lpfull}{\Lpc(\Omega, \B, m)}
\newcommand{\Linffull}{\Linfc(\Omega, \B, m)}
%%Lp après quotientage
\newcommand{\Li}{L^1}
\newcommand{\Ld}{L^2}
\newcommand{\Lp}{L^p}
\newcommand{\Linf}{L^{\infty}}
\newcommand{\Listd}{\Li(\Omega, m)} 
\newcommand{\Ldstd}{\Ld(\Omega, m)} 
\newcommand{\Lpstd}{\Lp(\Omega, m)} 

%Quelques opérateurs
\DeclareMathOperator{\codim}{codim}
\DeclareMathOperator{\rg}{rg}
\DeclareMathOperator{\tr}{tr}
\DeclareMathOperator{\vect}{Vect}
\DeclareMathOperator{\ima}{Im}
\DeclareMathOperator{\id}{Id}
\DeclareMathOperator{\sign}{sign}
\DeclareMathOperator{\ext}{ext}
\newcommand{\ind}{\mathbbm{1}}
\newcommand*{\spr}[2]{\left(\,#1\,\middle|\,#2\,\right)}
\newcommand{\interior}[1]{%
  {\kern0pt#1}^{\mathrm{o}}%
}
\DeclareMathOperator{\card}{card}
\DeclareMathOperator{\ppcm}{ppcm}

%Commandes de pure flemme
\newcommand{\veps}{\varepsilon}
\newcommand{\intab}{\int_{a}^{b}}
\newcommand{\intR}{\int_{-\infty}^{\infty}}
\newcommand{\intinf}{\int_{- \infty}^{+ \infty}}
\newcommand{\equi}{\Leftrightarrow}
\newcommand{\Kev}{$\K$-espace vectoriel }
\newcommand{\Kevs}{$\K$-espaces vectoriels }
\newcommand{\Rev}{$\R$-espace vectoriel }
\newcommand{\Revs}{$\R$-espaces vectoriels }
\newcommand{\Cev}{$\Cx$-espace vectoriel }
\newcommand{\Cevs}{$\Cx$-espaces vectoriels }
\newcommand{\evn}{(E, \norm{.})}
\newcommand{\interf}[2]{[#1,#2]}
\newcommand{\limb}{\lim\limits_}
\newcommand{\lebn}{\lambda_n}
\newcommand{\pp}{\text{~presque partout}}
\newcommand{\derivp}[2]{\frac{\partial #1}{\partial #2}}
\newcommand{\famiu}{(u_i)_{i \in I}}
\newcommand{\sev}{\text{~s.e.v.}}
\newcommand{\Mnp}{M_{n,p}(\K)}
\newcommand{\Mn}{M_{n}(\K)}
\newcommand{\tq}{\text{~t.q.~}}
\newcommand{\mq}{~montrer~que~}
\newcommand{\et}{\quad \text{et} \quad}
\newcommand{\ou}{\text{~ou~}}
\newcommand{\Kx}{\K[X]}


%Commandes de gars chelou
\renewcommand{\subset}{\subseteq}

%Commande restriction, ty duckduckgo
\def\restr#1#2{\mathchoice
              {\setbox1\hbox{${\displaystyle #1}_{\scriptstyle #2}$}
              \restraux{#1}{#2}}
              {\setbox1\hbox{${\textstyle #1}_{\scriptstyle #2}$}
              \restraux{#1}{#2}}
              {\setbox1\hbox{${\scriptstyle #1}_{\scriptscriptstyle #2}$}
              \restraux{#1}{#2}}
              {\setbox1\hbox{${\scriptscriptstyle #1}_{\scriptscriptstyle #2}$}
              \restraux{#1}{#2}}}
\def\restraux#1#2{{#1\,\smash{\vrule height .8\ht1 depth .85\dp1}}_{\,#2}} 

%Quelques normes
\DeclarePairedDelimiter\abs{\lvert}{\rvert}
\DeclarePairedDelimiter\labs{\bigl\lvert}{\bigr\rvert}
\DeclarePairedDelimiter\norm{\lVert}{\rVert}
\DeclarePairedDelimiterXPP\normi[1]{}\lVert\rVert{_1}{\ifblank{#1}{\:\cdot\:}{#1}}
\DeclarePairedDelimiterXPP\normd[1]{}\lVert\rVert{_2}{\ifblank{#1}{\:\cdot\:}{#1}}
\DeclarePairedDelimiterXPP\normp[1]{}\lVert\rVert{_p}{\ifblank{#1}{\:\cdot\:}{#1}}
\DeclarePairedDelimiterXPP\normq[1]{}\lVert\rVert{_q}{\ifblank{#1}{\:\cdot\:}{#1}}
\DeclarePairedDelimiterXPP\norminf[1]{}\lVert\rVert{_\infty}{\ifblank{#1}{\:\cdot\:}{#1}}

%Bracket en angle
\DeclarePairedDelimiter\angl{\langle}{\rangle}

%Set, parenthèses
\newcommand{\set}[1]{\{ #1 \}}
\newcommand{\bigset}[1]{\bigl\{ #1 \bigr\}}
\newcommand{\Bigset}[1]{\Bigl\{ #1 \Bigr\}}
\newcommand{\bigpar}[1]{\bigl( #1 \bigr)}
\newcommand{\Bigpar}[1]{\Bigl( #1 \Bigr)}

%Opitmisation des opérateurs de comparaison
\DeclarePairedDelimiter\tio{o (}{)}
\DeclarePairedDelimiter\gro{O (}{)}


\DeclareMathOperator{\fix}{Fix}
\DeclareMathOperator{\stab}{Stab}

\pagestyle{fancy}

\fancyhead[C]{}
\fancyhead[R]{}

\begin{document}

\underline{Source :} Berhuy - Algèbre (le grand combat) - page 61

\underline{Leçons :}
\begin{itemize}
\item 101 - Groupe opérant sur un ensemble. Exemples et applications.
\end{itemize}

On suppose que $E$ est un ensemble fini de cardinal $n$ et $G$ un groupe fini agissant sur $E$.

\begin{defn}[Fixateur]
Pour tout $g \in G$, on note :
\[ \fix(g) = \set{x \in E ; g \cdot x = x } , \]
le fixateur de $g$, i.e. l'ensemble des points de $E$ laissés fixes par l'action de $g$.
\end{defn}

\begin{thm}[Formule de Burnside]
En notant $\Omega$ l'ensemble des orbites de $E$ sous l'action de $G$, on a l'égalité :
\[ \card(\Omega) = \dfrac1{\abs{\Omega}} \sum_{g \in G} \abs{\fix(g)} . \]
\end{thm}

\begin{proof}
Notons :
\[ C = \set{ (g,x) \in G \times E ; g \cdot x = x } , \]
on va calculer le cardinal de $C$ de deux manières différentes.
On remarque d'abord que :
\[ C = \bigcup_{x \in E} \stab_G(x) \times \set{x} . \]
Cette union étant disjointe, on a :
\begin{align*}
\abs{C} & = \sum_{x \in E} \card(\stab_G(x) \times \set{x}) \\
& = \sum_{x \in E} \abs{ \stab_G(x) } \\
& = \sum_{x \in E} \dfrac{\abs{G}}{\abs{O_x}} 
 = \abs{G} \sum_{x \in E} \dfrac1{\abs{O_x}} \\
\end{align*}
Or, l'ensemble des orbites sous l'action de $G$ forme une partition de $E$. 
De plus, pour $x \in \omega \subset \Omega$ on a $O_x = \omega$, on peut donc écrire :
\begin{align*}
\abs{C} & = \abs{G} \sum_{\omega \in \Omega} \sum_{x \in \omega} \dfrac1{\abs{O_x}} \\
& = \abs{G} \sum_{\omega \in \Omega} \sum_{x \in \omega} \dfrac1{\abs{\omega}} \\
& = \abs{G} \sum_{\omega \in \Omega} \abs{\omega} \dfrac1{\abs{\omega}} \\
& = \abs{G} \abs{\Omega}
\end{align*}
D'autre part, on peut écrire :
\[ C = \bigcup_{g \in G} \set{g} \times \fix(g) , \]
et cette union étant aussi distincte on déduit :
\[ \abs{C} = \sum_{g \in G}  \abs{\fix(g)} , \]
ce qui donne directement la formule recherchée.
\end{proof}

Dans la suite on supposera que le groupe $G$ est inclus dans l'ensemble des permutations de $E$, 
de sorte que $G$ agisse naturellement sur $E$. 
Pour $x \in E$ et $g \in G$, on note $g(x)$ la permutation de $x$ par $g$ et on appelle $g$-orbite de $x$ l'ensemble des éléments
de $E$ obtenu par applications successives de $g$ à $x$.
On rappelle que l'ensemble de ces $g$-orbites forment une partition de $E$.

\begin{defn}[Coloriage]
On appelle coloriage de $E$ en au plus $q \in \Ne$ couleurs une application :
\[ f : E \to \set{1, \dots, q} . \]
On note $C(E, q)$ l'ensemble de ces coloriages. Pour $f \in C(E,q)$ on appelle type de $f$ le $q$-uplet :
\[ \xi(f) = (n_1, \dots, n_q) \subset \set{1, \dots, n}^q, \]
i.e. le nombre d'utilisation de chaque couleur, ou encore :
\[ \xi(f) = ( \abs{f^{-1}(1)}, \dots, \abs{f^{-1}(q)} )  . \]
On note $C(E,q,\xi)$ l'ensemble des coloriages de type $\xi \in \set{1, \dots, n}^q$.
\end{defn}

\begin{lem}
\label{lem:1}
L'application :
\[ \begin{array}{c c c}
G \times C(E,q) & \to & C(E,q) \\
(g,f) & \mapsto & g \cdot f ,
\end{array} \]
avec $g \cdot f : x \mapsto f(g^{-1}(x))$, définit une action de $G$ sur $C(E,q)$.
De plus, elle induit, par restriction, une action de $G$ sur $C(E,q,\xi)$.
\end{lem}

\begin{proof}
Pour $x \in E, f \in C(E,q)$, on a bien $(1_G \cdot f)(x) = f(1_G(x)) = f(x)$ soit $1_G \cdot f = f$. \\
Pour $g, g' \in G$, on a :
\begin{align*}
(g \cdot (g' \cdot f))(x) & = (g' \cdot f)(g^{-1}(x)) \\
& = f(g'^{-1} (g^{-1} (x))) \\
& = f((gg')^{-1}(x)) \\
& = (gg' \cdot f)(x)
\end{align*}
on a bien une action de $G$ sur $C(E,q)$. \\
On remarque de plus que le type d'un coloriage est invariant par action de $G$,
en effet pour $g \in G$ et $f \in C(E,q)$ on a :
\[ (g \cdot f)(x) = i \iff f(g^{-1}(x)) = x , \]
et donc : 
\[ x \in (g \cdot f)^{-1}(\set{i}) \iff g^{-1}(x) \in f^{-1}(\set{i}) \iff x \in g(f^{-1}(\set{i}) ). \]
Donc $(g \cdot f)^{-1}(\set{i}) = g(f^{-1}(\set{i}) )$ et comme $g \in S(E)$ est une bijection on déduit que :
\[ \abs{(g \cdot f)^{-1}(\set{i}) } = \abs{ f^{-1}(\set{i}) } , \]
ce qui revient à dire que l'action de $G$ préserve le type, et ce qui permet de restreindre l'action.
\end{proof}

\begin{defn}
On dit que deux coloriages $f, f' \in C(E,q)$ (resp. de $C(E,q, \xi)$) sont $G$-équivalents s'ils sont dans la même orbite 
sous l'action de $G$. On note $C_G(E,q)$ (resp. $C_G(E,q,\xi)$) l'ensemble des orbites des coloriages sous l'action de G
(resp. l'ensemble des orbites des coloriages de type $\xi$ sous l'action de $G$).
\end{defn}

\begin{rmq}
Cela correspond à l'idée intuitive du coloriage d'un objet : il est invariant, quelle que soit la façon dont on observe l'objet.
\end{rmq}

\begin{nota}
Soient $X_1, \dots, X_q, q$ indéterminées. Pour un type $\xi=(n_1, \dots, n_q)$, on note $X^{\xi}$ le monôme :
\[ X^{\xi} = X_1^{n_1} \cdots X_q^{n_q} . \]
\end{nota}

\begin{defn}[Inventaire]
On appelle inventaire de $E$ sous l'action de $G$ le polynôme :
\[ W_{E,G} = \sum_{\xi \in \set{1, \dots, n}^q} \abs{C_G(E,q,\xi)} X^{\xi} \in \Q[X_1, \dots, X_q] . \]
\end{defn}

\begin{lem}
\label{lem:2}
On a :
\[ W_{E,G} = \dfrac1{\abs{G}} \sum_{g \in G} \sum_{f \in C(E,q)(g)} X^{\xi(f)} , \]
en notant $C(E,q)(g) = \set{f \in C(E,q) ; g \cdot f = f}$ le fixateur de $g$ relativement à l'action de $G$ sur $C(E,q)$.
\end{lem}

\begin{proof}
Par application directe de la formule de Burnside, on a :
\[ \abs{C_G(E,q,\xi)} = \dfrac1{\abs{G}} \sum_{g \in G} \abs{C(E,q,\xi)(g)} , \]
avec $C(E,q, \xi)(g) = \set{f \in C(E,q, \xi) ; g \cdot f = f}$ le fixateur de $g$ relativement à l'action de $G$ sur $C(E,q, \xi)$.
En sommant sur les types on obtient :
\begin{align*}
\sum_{\xi \in \set{1, \dots, n}^q} \abs{C_G(E,q,\xi)} X^{\xi} & =  
\dfrac1{\abs{G}} \sum_{\xi \in \set{1, \dots, n}^q} \sum_{g \in G} \abs{C(E,q,\xi)(g)}  X^{\xi} \\
& = W_{E,G}
\end{align*}
Les variables $g$ et $\xi$ étant indépendantes, on peut écrire :
\[ W_{E,G} = \dfrac1{\abs{G}} \sum_{g \in G} \sum_{\xi \in \set{1, \dots, n}^q} \abs{C(E,q,\xi)(g)}  X^{\xi} . \]
Pour $g \in G$, on a :
\[ C(E,q)(g) = \bigcup_{\xi \in \set{1, \dots, n}^q} C(E,q,\xi)(g) , \]
cette union étant disjointe, ce qui permet d'écrire :
\begin{align*}
\sum_{f \in C(E,q)(g)} X^{\xi(f)} & = \sum_{\xi \in \set{1, \dots, n}^q} \sum_{f \in C(E,q,\xi)(g)} X^{\xi(f)} \\
& = \sum_{\xi \in \set{1, \dots, n}^q} \sum_{f \in C(E,q,\xi)(g)} X^{\xi} \\
& = \sum_{\xi \in \set{1, \dots, n}^q} \abs{C(E,q,\xi)(g)} X^{\xi} .
\end{align*} 
ce qui permet de conclure.
\end{proof}

\begin{lem}
\label{lem:3}
Soit $g \in G$. Un coloriage $f$ est élément de $C(E,q)(g)$ (i.e. est fixe sous l'action de $g$) si, et seulement si, $f$ est constant sur 
chaque $g$-orbite.
\end{lem}

\begin{proof}
Il est clair que si $g \cdot f = f$ alors $\forall x \in E, f(g^{-1}(x) = f(x)$ et on montre facilement que pour tout entier $p \in \N$,
et pour tout $x \in E, f(g^p(x)) = f(x)$. \\
Réciproquement, comme $g^{-1}(x)$ et $x$ sont dans le même $g$-orbite, il est clair que si $f$ est constant sur chaque $g$-orbite
alors il est fixe sous l'action de $G$.
\end{proof}

\begin{nota}
Fixons $g \in G$, il y a un nombre fini de $g$ orbite que l'on note $\omega_1, \dots, \omega_r$. 
Notant $Y= \set{1, \dots, r}$ on définit l'application :
\[ \begin{array}{c c c c}
\varphi : & C(Y,q) & \to & C(E,q)(g) \\
& g & \mapsto & \varphi(g) 
\end{array} \]
avec $\varphi(g)(x) = g(i)$, où $i$ est l'indice de $\omega_i$, la $g$-orbite de $x$. 
Par construction, $\varphi(g)$ est constant sur chaque $g$-orbite et est donc fixe sous l'action de $g$ donc $\varphi$ est bien définie.
Elle est de plus bijective car l'application $\psi$ définie sur $C(E,q)(g)$ par :
\[ \forall i \in \set{1, \dots, r}, \quad  \psi(f)(i) = f(x) , \]
avec $x$ un élément arbitraire de $\omega_i$ est bien définie. En effet, $f$ est fixe sous l'action de $g$ donc constante sur $\omega_i$
et donc l'image de $i$ par $\psi(f)$ ne dépend pas du représentant choisi. De plus, pour $g \in C(Y,q)$, on a :
\begin{align*}
\forall i \in \set{1, \dots, r}, \psi(\varphi(g))(i) & = \varphi(g)(x) \qquad (x \in \omega_i) \\
& = g(i) ,
\end{align*}
Autrement dit, $\psi(\varphi(g)) = g$ et de même $\varphi(\psi(f)) = f$ pour $f \in C(E,q)(g)$.
\end{nota}

\begin{lem}
\label{lem:4}
Soit $g \in G$ et soient $p_1, \dots, p_r$ les cardinaux des $g$-orbites $\omega_1, \dots, \omega_r$.  On a alors :
\[ \sum_{f \in C(E,q)(g)} X^{\xi(f)} = \prod_{i=1}^r(X_1^{p_i}+\dots+X_q^{p_i}) . \]
\end{lem}

\begin{proof}
Soit $f$ un coloriage, on a $X^{\xi(f)} = \prod_{x \in E} X_{f(x)}$. De plus, en utilisant la bijection $\varphi$, on a :
\[ \sum_{f \in C(E,q)(g)} X^{\xi(f)} = \sum_{g \in C(Y,q)} X^{\xi(\varphi(g))} = \sum_{g \in C(Y,q)} \prod_{x \in E} X_{\varphi(g)(x)} . \]
Comme les $g$-orbites forment une partition de $E$ et que $\varphi(g)$ est constante sur chacune de ces orbites, on a :
\[ \sum_{f \in C(E,q)(g)} X^{\xi(f)} = \sum_{g \in C(Y,q)} \prod_{i=1}^r X_g(i)^{p_i} . \]
En posant $g(i) = m_i$, on déduit :
\[ \sum_{f \in C(E,q)(g)} X^{\xi(f)} = \sum_{1 \leq m_1, \dots, m_r \leq q} \prod_{i=1}^r X_{m_i}^{p_i}
= \sum_{1 \leq m_1, \dots, m_r \leq q} X_{m_1}^{p_1} \cdots X_{m_r}^{p_r} . \] 
On reconnait alors le terme général du produit :
\[ \prod_{i=1}^r (X_1^{p_i} + \dots + X_q^{p_i}) , \]
ce qui permet de conclure.
\end{proof}

\begin{thm}[Polya]
L'inventaire de $E$ sous l'action de $G$ peut être exprimé par :
\[ W_{E,G} = \dfrac1{\abs{G}} \sum_{g \in G} \prod_{i=1}^n (X_1^i + \dots + X_q^i)^{e_i(g)} , \]
où $e_i(g)$ est le nombre de $g$-orbites à $i$ éléments.
\end{thm}

\begin{proof}
On remarque tout d'abord que les $g$-orbites sont toutes de cardinal compris entre 1 et $n$, autrement dit le nombre $e_i(g)$
est éventuellement non nul si, et seulement si, $i \in \set{1, \dots, n}$.
En reprenant la formule du lemme précédent, on peut écrire :
\[ \sum_{f \in C(E,q)(g)} X^{\xi(f)} = \prod_{i=1}^n (X_1^{p_i}+\dots+X_q^{p_i})^{e_i(g)}, \]
en regroupant les orbites de même cardinal (avec $e_i(g)$ éventuellement nul).
On conclut en substituant cette expression dans celle du lemme $\ref{lem:2}$.
\end{proof}

\begin{coro}
Le nombre $\abs{C_G(E,q)}$ de classes de $G$-équivalence de coloriages de $E$ en au plus $q$ couleurs est égal à :
\[ \abs{C_G(E,q)} = \dfrac1{\abs{G}} \sum_{g \in G} q^{r(g)} , \]
où $r(g)$ est le nombre de $g$-orbites.
\end{coro}

\begin{proof}
En utilisant l'union disjointe :
\[ C_G(E,q) = \bigcup_{\xi \in \set{1, \dots, n}^q} C_G(E,q, \xi) , \]
on déduit :
\[ \abs{C_G(E,q)} = \sum_{\xi \in \set{1, \dots, n}^q} \abs{C_G(E,q, \xi)} = W_{E,G}(1, \dots, 1) , \]
soit d'après le théorème de Polya :
\[ \abs{C_G(E,q)} = \dfrac1{\abs{G}} \sum_{g \in G} \prod_{i=1}^n q^{e_i(g)} = \dfrac1{\abs{G}} \sum_{g \in G} q^{\sum_{i=1}^n e_i(g)} , \]
et on conclut en remarquant que $\sum_{i=1}^n e_i(g) = r(g)$.
\end{proof}

\begin{appl}
Considérons un polygone régulier à $n$ sommets dont on veut colorer les arêtes avec au plus $q$ couleurs.
Les isométries possibles sont toutes les rotations d'angle $2k\pi/n$, avec $k \in \set{0, \dots, n-1}$, soit $\abs{G} = n$.
Le théorème de Lagrange nous dit que, pour tout élément $g \in G$, le cardinal du groupe cyclique $\angl{g}$ est un diviseur de $n$.
Donc, pour $g \in G$, la longueur des $g$-orbites est $d$, un diviseur de $n$, et il existe donc $n/d$ telles $g$ orbites.
Enfin, en utilisant l'égalité :
\[ n = \sum_{d | n} \varphi(d) , \]
avec $\varphi$ la fonctin indicatrice d'Euler.
On déduit qu'il existe exactement $\varphi(d)$ isométries dont l'orbite est de longueur $d$ (les isométries formant une partition de E).
Ce qui permet d'écrire :
\[ C = \dfrac1{n} \sum_{d | n} \varphi(d) q^{n/d} . \]
Par exemple, il existe 14 coloriages d'un hexagone régulier utilisant au plus 2 couleurs.
\end{appl}

\end{document}